%! Solving by Square Roots & Completing the Square — Unit 3, Lesson 3.3 — Homework (Answer Key)
\documentclass[10pt]{article}
\usepackage{algebra2-article}
\usepackage{algebra2-key}   % pulls in -boxes; do NOT also load -boxes
\newcommand{\TallMath}[1]{$\displaystyle #1\rule[-1.4em]{0pt}{3.2em}$}
\newcommand{\lgrid}[4]{%
  \draw[step=1, linegray!40, very thin] (#1,#3) grid (#2,#4);
  \draw[->, semithick, charcoal] (#1,0)--(#2,0) node[below right, font=\scriptsize]{$x$};
  \draw[->, semithick, charcoal] (0,#3)--(0,#4) node[left, font=\scriptsize]{$y$};}

\newcommand{\MethodA}{\ans{factoring}\hspace{0pt}}
\newcommand{\MethodB}{\ans{square roots}\hspace{0pt}}
\newcommand{\MethodC}{\ans{completing the square}\hspace{0pt}}

\begin{document}

\pageheader{Unit 3, Lesson 3.3}{Homework --- Answer Key}

\vspace{0.10in}

\begin{notesbox}{Practice}
Show your work. Isolate the square before you root, keep the $\pm$, and write radicals in simplest form.

\begin{enumerate}[leftmargin=1.9em, itemsep=11pt]
  \item \textbf{Solve by the Square Root Property.} Both roots, simplest radical form.
    \begin{center}
    \renewcommand{\arraystretch}{1.7}
    \begin{tabularx}{\linewidth}{@{}X X@{}}
    $x^2=81 \Rightarrow x=$ \ans{$\pm9$} & $x^2-10=0 \Rightarrow x=$ \ans{$\pm\sqrt{10}$} \\
    $4x^2-100=0 \Rightarrow x=$ \ans{$\pm5$} & $(x-1)^2=49 \Rightarrow x=$ \ans{$8,\ -6$} \\
    \multicolumn{2}{@{}l}{$(x+3)^2=8 \Rightarrow x=$ \ans{$-3\pm2\sqrt2$}} \\
    \end{tabularx}
    \end{center}

  \item \textbf{Complete the square to solve.} Show the completed square, then both roots.
    \begin{enumerate}[label=(\alph*), leftmargin=1.9em, itemsep=6pt]
      \item $x^2+6x+2=0$
        \begin{work}
          x^2+6x &= -2 \\
          x^2+6x+9 &= 7 \\
          (x+3)^2 &= 7 \\
          x &= -3\pm\sqrt{7}
        \end{work}
      \item $x^2-8x+5=0$
        \begin{work}
          x^2-8x &= -5 \\
          x^2-8x+16 &= 11 \\
          (x-4)^2 &= 11 \\
          x &= 4\pm\sqrt{11}
        \end{work}
      \item $x^2+2x-1=0$
        \begin{work}
          x^2+2x &= 1 \\
          x^2+2x+1 &= 2 \\
          (x+1)^2 &= 2 \\
          x &= -1\pm\sqrt{2}
        \end{work}
    \end{enumerate}

  \item \textbf{Rewrite in vertex form} $(x-h)^2+k$ by completing the square, and name the vertex.
    \begin{center}
    \renewcommand{\arraystretch}{1.6}
    \begin{tabularx}{\linewidth}{@{}l X X@{}}
    \textbf{Standard} & \textbf{Vertex form} & \textbf{Vertex} \\
    \hline
    $x^2+6x+1$ & \ans{$(x+3)^2-8$} & \ans{$(-3,-8)$} \\
    $x^2-4x-1$ & \ans{$(x-2)^2-5$} & \ans{$(2,-5)$} \\
    \end{tabularx}
    \end{center}

  \item \textbf{Choose a method, then solve.} Name it (\emph{factoring / square roots / completing the square}).
    \begin{enumerate}[label=(\alph*), leftmargin=1.9em, itemsep=6pt]
      \item $x^2-5x+6=0$ \hfill method: \MethodA
        \begin{work}
          (x-2)(x-3) &= 0 \\
          x &= 2,\ 3
        \end{work}
      \item $x^2=48$ \hfill method: \MethodB
        \begin{work}
          x &= \pm\sqrt{48} \\
            &= \pm4\sqrt{3}
        \end{work}
      \item $x^2+10x+18=0$ \hfill method: \MethodC
        \begin{work}
          x^2+10x &= -18 \\
          x^2+10x+25 &= 7 \\
          (x+5)^2 &= 7 \\
          x &= -5\pm\sqrt{7}
        \end{work}
    \end{enumerate}

  \item \textbf{Roots from a graph.} The graph of $f(x)=x^2-2x-4$ is shown; its roots are irrational.
    \begin{center}
    \begin{tikzpicture}[scale=0.5]
      \lgrid{-3}{5}{-6}{4}
      \begin{scope}\clip (-3,-6) rectangle (5,4);
        \draw[very thick, forest, domain=-2.0:4.0, samples=80, smooth] plot (\x,{\x*\x - 2*\x - 4});
      \end{scope}
      \fill[charcoal] (-1.236,0) circle (3pt);
      \fill[charcoal] (3.236,0) circle (3pt);
      \fill[charcoal] (1,-5) circle (3pt) node[below, font=\scriptsize]{$(1,-5)$};
    \end{tikzpicture}
    \end{center}
    Complete the square to find the \emph{exact} roots: $x^2-2x-4=0 \Rightarrow (x-1)^2=$ \ans{$5$},
    so $x=$ \ans{$1\pm\sqrt5$}. \; Do these match the two $x$-intercepts on the graph? \ans{yes}

  \item \textbf{Explain.} A classmate solves $x^2=25$ and writes $x=5$. What did they miss, and what is the
        complete answer? \ansline{They missed the negative root: $x=\pm5$. A number has two square roots, so $x=5$ \emph{or} $x=-5$.}
\end{enumerate}
\end{notesbox}

\begin{extensionbox}
\textbf{A matted photo.} A square photo of side $x$ inches gets a uniform $2$-inch mat all around, so the
whole framed square has side $(x+4)$ inches. The framed square's area is $121$ in$^2$.
\begin{enumerate}[label=(\alph*), leftmargin=1.8em, itemsep=5pt]
  \item Write an equation for the framed area using a square.
  \item Solve it by the Square Root Property.
  \item One root cannot be a side length. State the photo's side $x$ and explain why you rejected the other.
\end{enumerate}
\begin{work}
  (x+4)^2 &= 121 \\
  x+4 &= \pm11 \\
  x &= 7 \quad\text{or}\quad x = -15 \\
  x &= 7\text{ in} \quad\text{(reject } -15\text{: a side length can't be negative)}
\end{work}
\end{extensionbox}

\begin{spiralbox}
\textbf{Coming next --- Lesson 3.4: Complex Numbers.} Every method so far stops at the same wall: what is
$\sqrt{-9}$? Solving $x^2=-9$ or $x^2+2x+5=0$ leads to a \emph{negative} under the square root, so there is
\emph{no real} solution. Next you'll meet the imaginary unit $i=\sqrt{-1}$ and the \emph{complex numbers},
which finally let every quadratic be solved.
\end{spiralbox}

\end{document}
